\newcommand{\tipoRequerimento}{Isenção de ICMS}

\section{DOS FATOS E DOS DIREITOS}

\begin{enumerate}[resume]
    \item O art. 265, inciso XVII, do Regulamento do ICMS do Estado da Bahia (Decreto nº 13.780, de 16 de março de 2012) estabelece a isenção do valor do imposto sobre as operações com energia elétrica.

    \Citacao{
        Art. 265. São isentas do ICMS:

        [...]

        XLVII – as operações com energia elétrica quando destinada ao consumo em órgãos da administração pública municipal e fundações mantidas pelo poder público municipal e à iluminação pública;

        [...]

        XVII – as operações internas efetuadas por concessionário de serviço público de abastecimento de água (Conv. ICMS 98/89): 

        a)	nos fornecimentos de água natural canalizada: 

        1 – em ligações com consumo medido de até 30m3 mensais, por economia; 

        2 – em ligações não medidas, para consumo residencial; 

        3 –em ligações para consumo de órgãos públicos federais, estaduais e municipais; 

        b)	nos fornecimentos de água natural através de carro-pipa;
    }

    \item Dessa forma, as Unidades Consumidoras administradas pelo Serviço Público de água, esgoto e saneamento estão isentas da cobrança do ICMS.

    \item Foi constatado, por meio das faturas das UCs associadas à Prefeitura deste Município, que houve a cobrança de ICMS na fatura de energia elétrica da UC listada abaixo, sem a aplicação de qualquer isenção ou redução do imposto.

    \item Segue abaixo o número da Unidade Consumidora que necessita da correção tarifária, de responsabilidade do Município: \textbf{\unidadeConsumidora{}}.

    \item A REN nº 1000/2021 da ANEEL, em seu art. 323, determina as ações a serem adotadas em casos de cobrança errada decorrente de falhas atribuíveis à distribuidora. A referida resolução especifica que a restituição, quando houver cobranças excessivas, deve abranger todos os componentes da fatura, incluindo os tributos, conforme o trecho a seguir.
    
    \Citacao{
        Art. 323. A distribuidora, no caso de faturar valores incorretos, não apresentar fatura ou faturar sem utilizar a leitura do sistema de medição nos casos em que não haja previsão nesta Resolução, sem prejuízo das penalidades cabíveis, deve observar os seguintes procedimentos:

        [...]

        II - Faturamento a maior: devolver ao consumidor e demais usuários, até o 2º (segundo) ciclo de faturamento posterior à constatação, as quantias recebidas indevidamente nos últimos 60 (sessenta) ciclos de faturamento imediatamente anteriores à constatação.

        [...]

        § 2º No caso do inciso II do caput, a distribuidora deve devolver de acordo com as seguintes disposições: I - a quantia recebida indevidamente deve ser devolvida em dobro, independentemente de dolo ou culpa da distribuidora, salvo hipótese do §3º; 

        II - o valor do inciso I deste parágrafo deve ser atualizado pelo Índice Nacional de Preços ao Consumidor Amplo – IPCA; e 

        III - devem ser calculados e acrescidos os juros de mora à razão de 1\% (um por cento) ao mês pro rata die sobre o valor atualizado obtido do inciso II deste parágrafo.

        [...]

        § 4º A devolução em dobro prevista no § 2º é aplicável a todos os valores que compõem o faturamento, inclusive tributos, compensações, bandeiras tarifárias e cobranças de qualquer natureza.

        § 5º Aplica-se a devolução prevista no inciso II do caput no caso de cobranças adicionais decorrentes de erros da distribuidora que resultem em redução de benefícios tarifários ou aumentos tributários. (grifou-se)
    }

    \item Assim, destaca-se que a COELBA não pode apenas alegar ilegitimidade na cobrança, visto que a situação apresentada está respaldada pela REN nº 1000/2021, que impõe à concessionária a responsabilidade por erros que gerem acréscimos nos tributos. Portanto, a obrigação de devolver os valores recai sobre a distribuidora.

    \item A COELBA não pode se valer da REN nº 800/2017 como justificativa, uma vez que tal norma se aplica exclusivamente aos pedidos de reclassificação de unidades consumidoras. Referida resolução dispõe que, ausente solicitação no prazo de 90 (noventa) dias, o pedido deve ser tratado como nova ligação, afastando a devolução de valores apenas nesse contexto. O caso em análise, contudo, não trata de reclassificação, mas da restituição do ICMS indevidamente cobrado, em decorrência de erro da própria COELBA na aplicação do Regulamento do ICMS do Estado da Bahia.

    \item A reclamação não se vincula à REN nº 800/2017. O erro de enquadramento cometido pela COELBA não afasta o direito à restituição do ICMS indevidamente cobrado, uma vez que a Unidade Consumidora (UC) está vinculada ao CNPJ do Município, destinando-se ao consumo por órgãos da Administração Pública Municipal ou fundações por ela mantidas. O Regulamento do ICMS do Estado da Bahia é expresso ao isentar tais unidades da incidência de ICMS sobre energia elétrica.
\end{enumerate}

\section{DOS PEDIDOS}

\begin{enumerate}[resume]
    
\item Ante o exposto, requer-se:

\begin{enumerate}

    \item Proceda com a devolução de todos os valores cobrados a maior nos últimos 10 (dez) anos, conforme o Despacho nº 2.006 da ANEEL, de 08 de julho de 2024, observando o prazo prescricional de 10 (dez) anos, nos termos do art. 205 do Código Civil. Tal pedido encontra fundamento na sentença proferida em 29 de setembro de 2023 pela 19ª Vara Cível Federal de São Paulo, no âmbito da Ação Civil Pública Cível nº 5024153-93.2018.4.03.6100, bem como no Ofício nº 00710/2024/PFANEEL/PGF/AGU, de 08 de maio de 2024, e no Ofício nº 00937/2024/PFANEEL/PGF/AGU, de 12 de junho de 2024.

    \item Devolva os valores cobrados a maior em dobro, atualizados monetariamente pelo IPCA e acrescidos de juros de mora de 1\% (um por cento) ao mês, calculados \textit{pro rata die}, em conformidade com o estabelecido na Resolução Normativa nº 1.000/2021 da ANEEL.

    \item Proceda à análise e ao atendimento deste requerimento no prazo máximo de 10 (dez) dias úteis, informando o resultado da classificação tarifária, caso seja necessária a realização de vistoria, em conformidade com o estabelecido na Resolução Normativa nº 1.000/2021 da ANEEL.

    \item Proceda com a devolução do valor integral por meio de depósito bancário na conta corrente do Município, conforme estabelece o § 7º do art. 323 da Resolução Normativa nº 1.000/2021 da ANEEL, comunicando, por escrito, o resultado do deferimento desta reclamação por meio do endereço eletrônico \textit{abelgomescunha.adv@gmail.com}.
\end{enumerate}

\end{enumerate}